\section{Additional Interpretation and Tree-Search Controls}
\label{app:interpretation}

This appendix retains the detailed weight-transfer, near-optimality, model-mismatch, and tree-search analyses summarized in Section~\ref{sec:discussion}. The comparisons use the protocols and reference files stated locally. They describe the tested settings rather than general regime boundaries.

\subsection{The information-unit weight and tuning}
EFE fixes the exploration--exploitation weight at $w{=}1$ rather than eliminating it (Section~\ref{subsec:formal_rho_equiv}). The Pareto analysis (Figure~\ref{fig:pareto}) shows this weight sits inside the reward-maximizing tied bracket $w^*_{\text{ret}}$ on three of the five swept environments, while the success-maximizing $w^*_{\text{succ}}$ used for Planning+IG in the main tables lies at $10$--$50$. The gap between $w^*_{\text{ret}}$ and $w^*_{\text{succ}}$ has practical implications. $w{=}1$ is tuned to neither criterion. Its objective is reward plus information gain at weight one, and the fact that its measured reward sits with the reward-maximizing weights rather than with the success-maximizing ones on three of the five swept environments is an empirical finding of the sweep, not a design target. In safety-critical settings where near-certain accuracy is required, a higher weight may be appropriate, though such weights must still be tuned per-environment. EFE provides a principled starting point rather than a scale-free automatic calibration. Section~\ref{sec:pareto} quantifies the accuracy-for-reward trade, which on Tiger and Diagnosis buys at most $1.8$pp of accuracy for a reward cost that arrives entirely at the single grid step to $w{=}50$. The weight's dependence on the reward scale and on the bit-versus-nat convention, with the per-environment consequences of the latter, is set out in the reward-to-nats calibration discussion in Appendix~\ref{app:theory_setup}, and Section~\ref{sec:pareto} reports the measured sweeps behind those consequences.

A further complication for tuning is that adding planning depth to a weighted IG bonus amplifies the weight's effect, so the tuning problem becomes harder with depth. The same weight buys more sensing at depth two than at depth one, so the map from weight to sensing usage shifts with the horizon. A weight tuned at one depth need not carry over to another. One measurement supports this directly. At $w{=}100$ on Bandit under the canonical 5-seed protocol, Planning+IG at $H{=}2$ takes $12.36$ inspections against myopic Info Gain's $11.50$ (seed-level Welch $p{=}0.003$, \url{experiments/run_bandit_w100_depth_comparison.py}).

The requirement for a sufficient planning horizon also has a discount-factor form. In the discount sweep (Appendix~\ref{app:discount}, 5 seeds), the information-unit weight's advantage on the multi-observation environments appears at $\gamma \geq 0.99$ and not at the two lower values tested, and the transition is sharp. On Diagnosis ($N{=}4$), EFE's success advantage over Planning is significant at $\gamma \in \{0.99, 1.0\}$ and not statistically significant at $\gamma \in \{0.95, 0.90\}$ (Table~\ref{tab:discount}). The mechanism is that discounting truncates the effective planning horizon below the number of observations needed for confident diagnosis, erasing EFE's advantage. On Tiger (single observation), EFE and Planning differ significantly only at the heaviest discounting tested ($\gamma{=}0.90$) and are numerically identical at $\gamma \geq 0.95$. This is a meaningful practical limitation within the range swept. The sweep covers four discount factors on three environments, so it does not establish a general threshold, but it does show that applications whose time pressure mandates discounting as heavy as $\gamma{=}0.95$ on Diagnosis-like problems should not expect EFE's epistemic drive to pay off.

\subsection{Zero-shot weight transfer}
The tuning problem is not merely inconvenient. Whether a tuned weight survives transfer to a new environment depends on the tuning criterion. Table~\ref{tab:transfer} evaluates $w{=}1$, reward-tuned weights $w^*_{\mathrm{ret}}$, and success-tuned weights $w^*_{\mathrm{succ}}$ across environments. Success-tuned transfer harms reward at the top of the weight grid. $w^*_{\mathrm{succ}}{=}50$ is the success-tuned weight for both Tiger and Diagnosis, tuned within the Planning+IG agent class that the table evaluates. It yields $+4.25$ on Tiger, where it was tuned, and $-0.38$ on Testbed, where that same weight is carried to an environment it was not tuned for, versus $w{=}1$ at $+5.33$ and $+0.37$. Reward-tuned moderate weights ($w\in\{0.5,1\}$) transfer well within the shared reward scale. On Tiger and Diagnosis, $w{=}1$ is tied for best, and the tie is four-way, shared with $w{=}0.5$, Testbed's $w^*_{\mathrm{succ}}{=}10$, and Bandit's $w^*_{\mathrm{succ}}{=}20$. The reason is that the reward surface on each of those two environments is flat enough at moderate weights that reward-tuned and success-tuned weights from that band, whichever environment they were tuned on, all attain the same reward there. On Bandit, $w{=}1$ remains best outright without search. The honest conclusion is that moderate weights transfer robustly under a shared reward convention whatever their provenance, while success-tuned weights overfit once pushed to $w{=}50$, the top of the tested grid. In that moderate band, $w{=}1$ is the variational default, not a scale-free optimum.

\begin{table}[ht]
\centering
\caption{Zero-shot weight transfer (500 episodes $\times$ 5 seeds, reward reported as mean $\pm$ seed-level SE, \texttt{results\_\allowbreak transfer.csv}). Rows include $w{=}1$, reward-tuned weights ($w^*_{\mathrm{ret}}$), and success-tuned weights ($w^*_{\mathrm{succ}}$), the Planning+IG-class values recorded in \texttt{results\_\allowbreak supplementary\_\allowbreak tuned\_\allowbreak weights.csv}, matching the evaluated agent class. Success-tuned transfer harms reward at the top of the grid, while moderate weights transfer well. Bold entries mark the best reward per column (Tiger and Diagnosis show a four-way tie among $w \in \{0.5, 1, 10, 20\}$, bit-identical under the fixed pipeline, not merely within SE).}
\label{tab:transfer}
%% Numbers in this table are produced by experiments/run_transfer.py
\begin{tablebody}
\begin{tabular}{lcccc}
\toprule
Weight & Tiger & Diagnosis & Bandit & Testbed \\
\midrule
\shortstack[l]{$w{=}0.5$ ($w^*_{\mathrm{ret}}$)\\Diag./Testbed} & $\mathbf{+5.33 \pm 0.26}$ & $\mathbf{-1.26 \pm 0.34}$ & $+6.05 \pm 0.11$ & $\mathbf{+0.48 \pm 0.02}$ \\
\shortstack[l]{$w{=}1$ (EFE / $w^*_{\mathrm{ret}}$)\\Tiger/Bandit} & $\mathbf{+5.33 \pm 0.26}$ & $\mathbf{-1.26 \pm 0.34}$ & $\mathbf{+6.24 \pm 0.09}$ & $+0.37 \pm 0.01$ \\
\shortstack[l]{$w{=}10$ ($w^*_{\mathrm{succ}}$)\\Testbed} & $\mathbf{+5.33 \pm 0.26}$ & $\mathbf{-1.26 \pm 0.34}$ & $+5.25 \pm 0.06$ & $+0.01 \pm 0.01$ \\
\shortstack[l]{$w{=}20$ ($w^*_{\mathrm{succ}}$)\\Bandit} & $\mathbf{+5.33 \pm 0.26}$ & $\mathbf{-1.26 \pm 0.34}$ & $+5.05 \pm 0.07$ & $-0.18 \pm 0.01$ \\
\shortstack[l]{$w{=}50$ ($w^*_{\mathrm{succ}}$)\\Tiger/Diag.} & $+4.25 \pm 0.08$ & $-3.68 \pm 0.14$ & $+4.23 \pm 0.06$ & $-0.38 \pm 0.01$ \\
\bottomrule
\end{tabular}
\end{tablebody}
\end{table}

\paragraph{When is $w{=}1$ near-optimal?}
Having seen that $w{=}1$ transfers as a moderate default, we now ask when it is also near-optimal, first in closed form for the two-state case and then empirically across planning horizons. Proposition~\ref{prop:nearopt} formalizes lower and upper thresholds for the two-state case, written $w^*_{\mathrm{lo}}$ and $w^*_{\mathrm{hi}}$ in Table~\ref{tab:alpha_eta}. The lower threshold $w^*_{\mathrm{lo}}$ (the proposition's $w^*_{\mathrm{thresh}}$) is the weight below which the agent commits without observing at $H{=}1$, and it depends on the reward asymmetry $\alpha$, the informativeness $\eta$, and the observation accuracy $p$, all three defined in the proposition. The upper threshold $w^*_{\mathrm{hi}}$ (the proposition's $w^*_{\mathrm{over}}$) is the weight above which the agent, at the post-first-observation belief, buys a second observation whose instrumental value falls short of its cost, so it bounds over-observation. Part (a) of the proposition predicts that large reward asymmetry drives the lower threshold below zero, which leaves $w{=}1$ inside the interval whenever the upper threshold exceeds $1$. Appendix~\ref{app:theory_thresholds} compares Table~\ref{tab:alpha_eta} with that prediction. The high-asymmetry environments ($\alpha \geq 5$, namely Tiger, Diagnosis, and Tileworld) place $w{=}1$ interior to the interval, on the genuinely two-state row and, more heuristically, on the two-state-reduced rows, while Testbed places $w{=}1$ near the upper threshold, consistent with over-exploration.

The proposition's two parts live at two horizons, Part (a)'s lower threshold at $H{=}1$ and Part (b)'s near-optimality interval at $H{=}2$. To ask what happens to the near-optimality of $w{=}1$ as the planning horizon varies, including beyond those two cases, Appendix~\ref{app:nearopt_horizon} reports a Monte Carlo study over 100 randomly generated two-state environments at $H \in \{1, 2, 3\}$ (Figure~\ref{fig:nearopt_horizon}). This study uses a different, empirical criterion of near-optimality from the closed-form interval above. The interval is stated for the two-state case at $H{=}2$ only, whereas a reward-gap criterion, which asks how much reward $w{=}1$ gives up relative to the best weight found by grid search, can be applied at any horizon, and the study counts $w{=}1$ as near-optimal when that gap is within $\max(5\%, 0.5)$ of the best reward. The near-optimality rate increases from 30\% at $H{=}1$ to 58\% at $H{=}2$ and 79\% at $H{=}3$, consistent with multi-step planning widening the near-optimality basin, the region of environment parameters where $w{=}1$ passes this test. High-asymmetry environments ($\alpha \geq 10$) follow the same trend, rising from 26\% to 53\% to 79\%. The failure rate, the fraction of environments where $w{=}1$ does not pass this test, is the more useful number for a practitioner and we state it directly. $w{=}1$ fails this near-optimality test in about $21\%$ of the random family even at $H{=}3$, and in $70\%$ of it at $H{=}1$. The criterion is also lenient by construction, since the absolute floor of $0.5$ reward units does most of the work on environments whose reward range is small.

The widening with horizon of the basin in which $w{=}1$ passes this test is qualitatively consistent with what we observed on Bandit. Bandit's asymmetry is a contrast between two positive arm rewards (correct-arm reward $+10$ vs.\ small-reward-arm $+1$) rather than the two-state $\alpha{=}|R^-|/R^+$ ratio this section's formula assumes. The formula therefore does not apply to it directly and we do not report a numeric $\alpha$ for it. Bandit, read loosely as a low-asymmetry case relative to the high-asymmetry environments above, is where the one-step analysis would lead one to expect the least from $w{=}1$, although we do not compute its threshold. Multi-step EFE at $H{=}2$ achieves $86.9\%$ success on it (Table~\ref{tab:main}), consistent with the aggregate direction of the random-environment study, where the fraction of environments in which $w{=}1$ meets that study's reward-gap near-optimality criterion rises from 30\% at $H{=}1$ to 58\% at $H{=}2$.

Many real-world decision problems have $\alpha \gg 1$. In medical diagnosis, fault detection, and security screening, a missed condition costs far more than another test. This is the regime where the interval analysis is most favorable to $w{=}1$, and the multi-step study shows the near-optimality basin widening with planning depth. The practical consequence is that in these domains the information-unit weight is a strong untuned starting default under the shared reward convention, not a guarantee of near-optimality. When the default's induced sensing differs from a stated expected-usage target, the budgeted machinery of Section~\ref{sec:budget} selects, within the single-price receding-horizon family, the weight or endpoint mixture whose expected usage meets that target.

\subsection{Robustness to model misspecification}
Our main results assume exact knowledge of the generative model. Appendix~\ref{app:misspec} investigates robustness when the agent's believed observation accuracy differs from the true value, over mismatches from $-0.15$ to $+0.10$, where a negative mismatch means the agent believes the sensor to be less accurate than it is and a positive mismatch means the reverse. One question the analysis can answer is whether EFE's epistemic term buffers a misspecified sensor in a way that reward-only Planning does not. On Tiger, success stays close to its zero-mismatch value at every tested level, and EFE and Planning respond identically to every tested mismatch (Table~\ref{tab:misspec-tiger}), so there is no measurable EFE-specific buffering effect on this environment. On Diagnosis, degradation is larger but graceful, and both agents' behavior steps between two plateaus rather than degrading smoothly with the mismatch magnitude, a higher one with more tests and higher success and a lower one with fewer tests and lower success. Changing the mismatch moves an agent between the plateaus rather than along a gradient (Table~\ref{tab:misspec-diag}). Overestimating sensor accuracy (positive mismatch) is more harmful than underestimating it, because the agent commits prematurely on insufficient evidence.

\subsection{Approximate planning with MCTS-EFE}
At each depth the tree branches over $K$ observation actions and then, per action, over its $|\mathcal{O}|$ observation outcomes, so the exact tree search cost is $\mathcal{O}((K \cdot |\mathcal{O}|)^H)$, limiting practical horizons to $H{=}2$--$3$ with $K \geq 3$. To push beyond this, we implement MCTS-EFE, a UCB1 tree search whose design choices we state in turn. The tree branches over observe and commit actions and, after each observe action, enumerates its $|\mathcal{O}|$ observation outcomes as child nodes. It adds the exact per-node information gain to the in-tree reward of each observe action, so that within the tree an observation is credited for what it reveals as well as charged for what it costs. Values propagate by max-backup, each node taking the maximum over its children's values rather than the average of the returns sampled through it. Leaves are scored by an EFE-greedy rollout, which at each step takes the observation with the best one-step EFE value and credits the better of committing now and continuing to observe. The UCB1 exploration constant is left at its default, the commit reward range (110 on Tiger). That default is a static analog of the $R_{\mathrm{hi}} - R_{\mathrm{lo}}$ rule of \citet{silver2010}, which sizes the constant to the spread between the highest and lowest returns seen in preliminary sample runs of the search, whereas the default here takes the spread once from the commit reward table and never revisits it. On Tiger at $H{=}10$, MCTS-EFE with 500 simulations achieves $97.7\% \pm 0.5$pp success ($23$--$28$ ms per episode across the three horizons tested, \texttt{results\_\allowbreak mcts\_\allowbreak efe.csv}), compared to $88.8\% \pm 0.2$pp for POMCP at the same simulation count ($15$--$19$ ms per episode across horizons, seed-level SE, seed-level Welch $p{=}5.2{\times}10^{-6}$). The comparison is simulation-matched, not compute-matched, since MCTS-EFE's leaf evaluations are more expensive per simulation. In this configuration, EFE's closed-form information valuation, acting as the in-tree observe reward and inside the leaf rollouts, outperforms POMCP's semi-informed rollouts (Appendix~\ref{app:pomcp}).

The two solvers differ in more than the leaf, so an ablation narrows where that advantage comes from. Two of the components that distinguish this search from a UCT tree (Upper Confidence bounds applied to Trees, the standard UCB1 tree search) can be turned off independently, the in-tree information gain and the max-backup, and Table~\ref{tab:mcts-ablation} runs all four combinations on Tiger, Diagnosis, and Tileworld-$6{\times}6$. Neither component shows a detectable effect on success or reward at five seeds after Holm correction (Appendix~\ref{app:pomcp}). What remains to carry the contribution is the EFE-greedy leaf rollout together with the exact belief representation and the exact commit-child values, which the ablation holds fixed. A commit child is the node reached by a commit action, and the search sets its value to the belief-expected commit reward computed directly from the belief instead of estimating it from sampled returns.

A natural objection is that the POMCP gap reflects an untuned exploration constant rather than the solver. The observe-then-commit POMCP runs with a single untuned UCB1 exploration constant per battery, and Appendix~\ref{app:pomcp} sweeps that constant on Tiger, Diagnosis, and Tileworld-$6{\times}6$ (Table~\ref{tab:pomcp-sweep}). The sweep changes the Tiger reading substantially. At its best swept constant POMCP reaches $96.3\% \pm 0.6$pp success, and we cannot distinguish that configuration from MCTS-EFE at its default constant on either metric, that is, we fail to reject the null hypothesis of no difference (seed-level Welch $p{=}0.09$ and $p{=}0.54$). The $88.8\%$ figure above therefore measures an untuned constant rather than the solver. The sweep also shows MCTS-EFE's default is not its own best, so the appendix compares the two solvers with each at a constant selected on tuning seeds disjoint from every evaluation seed, under a rule recorded before the selection was run. At those constants, MCTS-EFE stays ahead on Tiger success ($98.2\% \pm 0.3$pp against $96.3\% \pm 0.6$pp, $p{=}0.024$, surviving Holm correction within the appendix's six-test family) with reward not distinguishable ($p{=}0.16$), so tuning narrows the Tiger gap sharply without removing it. Because the canonical runs had been inspected before the rule was written, that narrow Tiger result is retrospective evidence rather than a prospective confirmation, and a confirmatory claim would need fresh evaluation seeds that no one has looked at. The Diagnosis and Tileworld differences are large enough that the same caveat does not change their reading. Where the agent must choose among tests, tuning the constant leaves most of the gap in place. POMCP's selected configuration reaches $71.4\% \pm 1.5$pp on Diagnosis and $19.7\% \pm 1.7$pp on Tileworld-$6{\times}6$ against MCTS-EFE's $94.7\% \pm 0.4$pp and $95.4\% \pm 0.6$pp, every comparison significant after Holm correction.

A second objection is that the gap reflects POMCP's rollout policy. The comparison already avoids a pure-random-rollout strawman, since POMCP's semi-informed rollouts sample observations uniformly but end in a belief-optimal commit (Appendix~\ref{app:pomcp}). Information-gain rollouts do not close the Diagnosis and Tileworld gap either. At $c{=}5$ they put Diagnosis at $71.4\%$ against the uniform-rollout POMCP's $68.2$ at the same constant and Tileworld-$6{\times}6$ at $19.1\%$ against $11.3$, neither difference surviving Holm correction (seed-level Welch $p{=}0.25$ and $p{=}0.023$). On Tiger they change nothing at all, since a single observation action makes the informed and uniform rollout policies identical. Richer domain knowledge or a learned value function in the rollout could narrow that gap.

Taken together, the sweep and the rollout check show that the remaining gap is not a matter of the exploration constant or of the rollout policy. Where sensing is one repeated action and only its stopping time is in question, as on Tiger, tuning the constant recovers most of the difference. Where the agent must choose which of several tests to run, neither control removes more than a small part of the gap, which stays above twenty points on Diagnosis and above seventy points on Tileworld-$6{\times}6$ at every swept configuration. What the evidence supports as the explanation is MCTS-EFE's joint configuration of EFE-greedy leaf evaluation over exact beliefs and exact commit-child values, since those three components were held fixed together in every run of the sweep, the rollout check, and the ablation. We therefore do not attribute the gap to directed observation selection in isolation. Separating the leaf evaluation from the exact belief and commit values would need a further matched-compute comparison that we have not run. The claim is not that EFE-based planning exceeds state-of-the-art POMDP solvers, but that the EFE objective supplies a principled, zero-tuning in-tree reward and leaf evaluation in place of hand-designed rollouts.

The approach also scales to multiple observation actions. For single-observation environments such as Tiger, the observation-outcome enumeration (2 outcomes per action) keeps expansion costs low, and the question is whether the same holds once the agent must choose among several tests. On Diagnosis ($N{=}4$, $K{=}2$ tests), MCTS-EFE at $H{=}5$ achieves $94.2\% \pm 0.5$pp success compared to POMCP's $68.7\% \pm 1.8$pp at the same simulation count (200 simulations, seed-level SE, seed-level Welch $p{=}7.2{\times}10^{-5}$, \texttt{results\_\allowbreak mcts\_\allowbreak efe.csv}). To exercise the approach on a larger multi-observation environment, we evaluate on Tileworld $6{\times}6$ ($|S|{=}36$, $K{=}6$ scans). MCTS-EFE saturates at $95.4\% \pm 0.6$pp success across both simulation budgets tested (200 and 500, seed-level SE), well above the $71.9\% \pm 1.5$pp of Exact-EFE (the full recursive EFE agent, called EFE in the earlier tables, evaluated here within the same MCTS battery), but at a reward cost ($-25.02$ vs.\ Exact-EFE's $-21.62$) from scanning substantially more (32.3 vs.\ 14.8 observations on average). MCTS's sampled rollouts do not value information as precisely as the exact recursion, so it trades reward for a large success-rate gain by scanning more thoroughly. On Tiger and Diagnosis there is no such trade, since MCTS-EFE sits slightly below Exact-EFE on success and further below it on reward ($97.7\%$ against $99.9\%$ and $+3.73$ against $+5.65$ on Tiger, $94.2\%$ against $97.1\%$ and $-3.26$ against $-1.44$ on Diagnosis). POMCP collapses on this larger observation menu regardless of simulation budget ($10.7\% \pm 1.1$pp at 200 simulations, $5.6\% \pm 1.1$pp at 500, seed-level SE). MCTS-EFE's $85$--$90$pp success-rate gap over POMCP (seed-level Welch $p<10^{-9}$ at both budgets) therefore shows that the gap persists at this scale, even though the specific reward/success trade-off MCTS-EFE strikes here differs from the smaller environments. Semi-informed rollouts still cannot identify which of the 6 scans to perform. For larger observation spaces, more efficient tree policies, such as progressive widening or double progressive widening \citep{benchetrit2025}, techniques that limit how many children a node expands as a function of its visit count, would further improve scalability. We leave this to future work.
